% ============================================================
% preamble.tex — 计算化学方法讲义 公共导言区
% 文档类在 main.tex 中：ctexart[UTF8,11pt,a4paper,fontset=fandol]
% 本文件只负责宏包、颜色、宏与 tcolorbox 环境。
% ============================================================

% ---------- 页面与排版 ----------
\usepackage[a4paper,margin=2.5cm,headheight=15pt]{geometry}
\usepackage{amsmath,amssymb,mathtools,bm}
\usepackage{graphicx}
\usepackage{xcolor}

% ---------- 图与表 ----------
\usepackage{tikz}
\usetikzlibrary{arrows.meta,positioning,shapes.geometric,calc,fit,backgrounds}
\usepackage{pgfplots}
\pgfplotsset{compat=1.17}
\usepackage{booktabs}
\usepackage{longtable}
\usepackage{array}
\usepackage{caption}
\usepackage{subcaption}
\usepackage{enumitem}

% ---------- 代码 ----------
\usepackage{listings}

% ---------- 文字与页眉 ----------
\usepackage{microtype}
\usepackage{fancyhdr}

% ---------- tcolorbox（必须在 hyperref 之前） ----------
\usepackage[most]{tcolorbox}

% ============================================================
% Okabe–Ito 色板（全书唯一颜色来源；章节文件不得写颜色字面量）
% ============================================================
\definecolor{oiBlue}{HTML}{0072B2}
\definecolor{oiOrange}{HTML}{E69F00}
\definecolor{oiGreen}{HTML}{009E73}
\definecolor{oiVermillion}{HTML}{D55E00}
\definecolor{oiPurple}{HTML}{CC79A7}
\definecolor{oiSky}{HTML}{56B4E9}

% ---------- 引用与链接（hyperref 最后加载） ----------
\usepackage[numbers,sort&compress]{natbib}
\usepackage[unicode,colorlinks=true,bookmarksnumbered=true,
            linkcolor=oiBlue,citecolor=oiGreen,urlcolor=oiVermillion]{hyperref}

% ============================================================
% 数学宏（CONTRACT.md §3，章节作者只使用、不定义）
% ============================================================
\newcommand{\dd}{\mathrm{d}}
\newcommand{\ee}{\mathrm{e}}
\newcommand{\ii}{\mathrm{i}}
\newcommand{\HF}{\mathrm{HF}}
\newcommand{\xc}{\mathrm{xc}}
\newcommand{\ket}[1]{\lvert #1\rangle}
\newcommand{\bra}[1]{\langle #1\rvert}
\newcommand{\op}[1]{\hat{#1}}

% \braket{a|b}      -> \langle a \mid b \rangle
% \bracket{a|O|b}   -> \langle a \mid O \mid b \rangle
% 以 | 为定界符递归切分（1 个、2 个或更多竖线均可）；
% 每段之间排 \mid，绝不原样输出 '|'。
\makeatletter
\def\braket@nil{\@nil}
\def\braket#1{\begingroup\langle\braket@loop#1|\braket@nil\rangle\endgroup}
\def\braket@loop#1|#2\braket@nil{%
  #1%
  \if\relax\detokenize{#2}\relax\else\mid\braket@loop#2\braket@nil\fi}
\def\bracket#1{\begingroup\langle\bracket@loop#1|\braket@nil\rangle\endgroup}
\def\bracket@loop#1|#2\braket@nil{%
  #1%
  \if\relax\detokenize{#2}\relax\else\mid\bracket@loop#2\braket@nil\fi}
\makeatother

% ============================================================
% 四个可断页 tcolorbox 环境（均接受一个可选参数）
% ============================================================
\newtcolorbox{problemchain}[1][]{
  enhanced, breakable,
  title={问题链},
  colframe=oiBlue, colback=oiBlue!6!white, colbacktitle=oiBlue, coltitle=white,
  boxrule=0.6pt, arc=2pt, left=6pt, right=6pt, top=4pt, bottom=4pt,
  fonttitle=\bfseries, #1}

\newtcolorbox{methodcard}[1][]{
  enhanced, breakable,
  title={方法卡片},
  colframe=oiGreen, colback=oiGreen!6!white, colbacktitle=oiGreen, coltitle=white,
  boxrule=0.6pt, arc=2pt, left=6pt, right=6pt, top=4pt, bottom=4pt,
  fonttitle=\bfseries, #1}

\newtcolorbox{ch4box}[1][]{
  enhanced, breakable,
  title={CH$_4$ 算例（def2-TZVPP）},
  colframe=oiOrange, colback=oiOrange!6!white, colbacktitle=oiOrange, coltitle=black,
  boxrule=0.6pt, arc=2pt, left=6pt, right=6pt, top=4pt, bottom=4pt,
  fonttitle=\bfseries, #1}

\newtcolorbox{pitfall}[1][]{
  enhanced, breakable,
  title={常见陷阱},
  colframe=oiVermillion, colback=oiVermillion!6!white, colbacktitle=oiVermillion, coltitle=white,
  boxrule=0.6pt, arc=2pt, left=6pt, right=6pt, top=4pt, bottom=4pt,
  fonttitle=\bfseries, #1}

% ============================================================
% 编号、目录、页眉、参考文献
% ============================================================
\numberwithin{equation}{section}
\numberwithin{figure}{section}
\numberwithin{table}{section}
\allowdisplaybreaks
\setcounter{tocdepth}{2}
\setcounter{secnumdepth}{2}
\renewcommand{\refname}{参考文献}

\pagestyle{fancy}
\fancyhf{}
\fancyhead[L]{\small\nouppercase{\leftmark}}
\fancyhead[R]{\small 计算化学方法讲义}
\fancyfoot[C]{\thepage}
\renewcommand{\headrulewidth}{0.4pt}
\renewcommand{\footrulewidth}{0pt}

% 章标题自动生成“第 N 章”；章节文件只写 \section{<标题>}
\ctexset{section/name={第,章}}

% ============================================================
% ORCA 输出 listings 风格
% ============================================================
\lstdefinestyle{orca}{
  basicstyle=\ttfamily\small,
  breaklines=true,
  breakatwhitespace=false,
  frame=single,
  framerule=0.4pt,
  rulecolor=\color{oiBlue!45!white},
  backgroundcolor=\color{oiBlue!3!white},
  columns=fullflexible,
  keepspaces=true,
  showstringspaces=false,
  numbers=none,
  commentstyle=\color{oiGreen},
  keywordstyle=\color{oiBlue}\bfseries,
  stringstyle=\color{oiVermillion},
  xleftmargin=0.5em,
  xrightmargin=0.5em,
  aboveskip=0.5em,
  belowskip=0.5em
}
\lstset{style=orca}

% 列表间距
\setlist{itemsep=2pt,topsep=3pt}
